\documentclass[]{book}
\usepackage{lmodern}
\usepackage{amssymb,amsmath}
\usepackage{ifxetex,ifluatex}
\usepackage{fixltx2e} % provides \textsubscript
\ifnum 0\ifxetex 1\fi\ifluatex 1\fi=0 % if pdftex
  \usepackage[T1]{fontenc}
  \usepackage[utf8]{inputenc}
\else % if luatex or xelatex
  \ifxetex
    \usepackage{mathspec}
  \else
    \usepackage{fontspec}
  \fi
  \defaultfontfeatures{Ligatures=TeX,Scale=MatchLowercase}
\fi
% use upquote if available, for straight quotes in verbatim environments
\IfFileExists{upquote.sty}{\usepackage{upquote}}{}
% use microtype if available
\IfFileExists{microtype.sty}{%
\usepackage{microtype}
\UseMicrotypeSet[protrusion]{basicmath} % disable protrusion for tt fonts
}{}
\usepackage{hyperref}
\hypersetup{unicode=true,
            pdftitle={Urban Transportation Planning},
            pdfauthor={Gregory Macfarlane, PhD, PE},
            pdfborder={0 0 0},
            breaklinks=true}
\urlstyle{same}  % don't use monospace font for urls
\usepackage{natbib}
\bibliographystyle{apalike}
\usepackage{color}
\usepackage{fancyvrb}
\newcommand{\VerbBar}{|}
\newcommand{\VERB}{\Verb[commandchars=\\\{\}]}
\DefineVerbatimEnvironment{Highlighting}{Verbatim}{commandchars=\\\{\}}
% Add ',fontsize=\small' for more characters per line
\usepackage{framed}
\definecolor{shadecolor}{RGB}{248,248,248}
\newenvironment{Shaded}{\begin{snugshade}}{\end{snugshade}}
\newcommand{\AlertTok}[1]{\textcolor[rgb]{0.94,0.16,0.16}{#1}}
\newcommand{\AnnotationTok}[1]{\textcolor[rgb]{0.56,0.35,0.01}{\textbf{\textit{#1}}}}
\newcommand{\AttributeTok}[1]{\textcolor[rgb]{0.77,0.63,0.00}{#1}}
\newcommand{\BaseNTok}[1]{\textcolor[rgb]{0.00,0.00,0.81}{#1}}
\newcommand{\BuiltInTok}[1]{#1}
\newcommand{\CharTok}[1]{\textcolor[rgb]{0.31,0.60,0.02}{#1}}
\newcommand{\CommentTok}[1]{\textcolor[rgb]{0.56,0.35,0.01}{\textit{#1}}}
\newcommand{\CommentVarTok}[1]{\textcolor[rgb]{0.56,0.35,0.01}{\textbf{\textit{#1}}}}
\newcommand{\ConstantTok}[1]{\textcolor[rgb]{0.00,0.00,0.00}{#1}}
\newcommand{\ControlFlowTok}[1]{\textcolor[rgb]{0.13,0.29,0.53}{\textbf{#1}}}
\newcommand{\DataTypeTok}[1]{\textcolor[rgb]{0.13,0.29,0.53}{#1}}
\newcommand{\DecValTok}[1]{\textcolor[rgb]{0.00,0.00,0.81}{#1}}
\newcommand{\DocumentationTok}[1]{\textcolor[rgb]{0.56,0.35,0.01}{\textbf{\textit{#1}}}}
\newcommand{\ErrorTok}[1]{\textcolor[rgb]{0.64,0.00,0.00}{\textbf{#1}}}
\newcommand{\ExtensionTok}[1]{#1}
\newcommand{\FloatTok}[1]{\textcolor[rgb]{0.00,0.00,0.81}{#1}}
\newcommand{\FunctionTok}[1]{\textcolor[rgb]{0.00,0.00,0.00}{#1}}
\newcommand{\ImportTok}[1]{#1}
\newcommand{\InformationTok}[1]{\textcolor[rgb]{0.56,0.35,0.01}{\textbf{\textit{#1}}}}
\newcommand{\KeywordTok}[1]{\textcolor[rgb]{0.13,0.29,0.53}{\textbf{#1}}}
\newcommand{\NormalTok}[1]{#1}
\newcommand{\OperatorTok}[1]{\textcolor[rgb]{0.81,0.36,0.00}{\textbf{#1}}}
\newcommand{\OtherTok}[1]{\textcolor[rgb]{0.56,0.35,0.01}{#1}}
\newcommand{\PreprocessorTok}[1]{\textcolor[rgb]{0.56,0.35,0.01}{\textit{#1}}}
\newcommand{\RegionMarkerTok}[1]{#1}
\newcommand{\SpecialCharTok}[1]{\textcolor[rgb]{0.00,0.00,0.00}{#1}}
\newcommand{\SpecialStringTok}[1]{\textcolor[rgb]{0.31,0.60,0.02}{#1}}
\newcommand{\StringTok}[1]{\textcolor[rgb]{0.31,0.60,0.02}{#1}}
\newcommand{\VariableTok}[1]{\textcolor[rgb]{0.00,0.00,0.00}{#1}}
\newcommand{\VerbatimStringTok}[1]{\textcolor[rgb]{0.31,0.60,0.02}{#1}}
\newcommand{\WarningTok}[1]{\textcolor[rgb]{0.56,0.35,0.01}{\textbf{\textit{#1}}}}
\usepackage{longtable,booktabs}
\usepackage{graphicx,grffile}
\makeatletter
\def\maxwidth{\ifdim\Gin@nat@width>\linewidth\linewidth\else\Gin@nat@width\fi}
\def\maxheight{\ifdim\Gin@nat@height>\textheight\textheight\else\Gin@nat@height\fi}
\makeatother
% Scale images if necessary, so that they will not overflow the page
% margins by default, and it is still possible to overwrite the defaults
% using explicit options in \includegraphics[width, height, ...]{}
\setkeys{Gin}{width=\maxwidth,height=\maxheight,keepaspectratio}
\IfFileExists{parskip.sty}{%
\usepackage{parskip}
}{% else
\setlength{\parindent}{0pt}
\setlength{\parskip}{6pt plus 2pt minus 1pt}
}
\setlength{\emergencystretch}{3em}  % prevent overfull lines
\providecommand{\tightlist}{%
  \setlength{\itemsep}{0pt}\setlength{\parskip}{0pt}}
\setcounter{secnumdepth}{5}
% Redefines (sub)paragraphs to behave more like sections
\ifx\paragraph\undefined\else
\let\oldparagraph\paragraph
\renewcommand{\paragraph}[1]{\oldparagraph{#1}\mbox{}}
\fi
\ifx\subparagraph\undefined\else
\let\oldsubparagraph\subparagraph
\renewcommand{\subparagraph}[1]{\oldsubparagraph{#1}\mbox{}}
\fi

%%% Use protect on footnotes to avoid problems with footnotes in titles
\let\rmarkdownfootnote\footnote%
\def\footnote{\protect\rmarkdownfootnote}

%%% Change title format to be more compact
\usepackage{titling}

% Create subtitle command for use in maketitle
\providecommand{\subtitle}[1]{
  \posttitle{
    \begin{center}\large#1\end{center}
    }
}

\setlength{\droptitle}{-2em}

  \title{Urban Transportation Planning}
    \pretitle{\vspace{\droptitle}\centering\huge}
  \posttitle{\par}
    \author{Gregory Macfarlane, PhD, PE}
    \preauthor{\centering\large\emph}
  \postauthor{\par}
      \predate{\centering\large\emph}
  \postdate{\par}
    \date{2020-04-30}

\usepackage{booktabs}

\begin{document}
\maketitle

{
\setcounter{tocdepth}{1}
\tableofcontents
}
\hypertarget{syllabus}{%
\chapter*{Foreword}\label{syllabus}}
\addcontentsline{toc}{chapter}{Foreword}

This book contains course notes and assignments for a senior / graduate class in
transportation planning and elementary travel modeling. A description for this course
is:

\begin{quote}
An advanced course in urban transportation planning. Urban transportation as
the outcome of an economic system, details and techniques for four-step travel
model development, applications of travel models within a legal and regulatory
context.
\end{quote}

The book is organized into six units:

\begin{enumerate}
\def\labelenumi{\arabic{enumi}.}
\tightlist
\item
  \protect\hyperlink{chap-blocks}{Building Blocks}
\item
  \protect\hyperlink{chap-tripgen}{Trip Generation}
\item
  \protect\hyperlink{chap-distribution}{Trip Distribution}
\item
  \protect\hyperlink{chap-modechoice}{Mode and Destination Choice}
\item
  \protect\hyperlink{chap-assignment}{Network Assignment and Validation}
\item
  \protect\hyperlink{chap-process}{The Planning Process}
\end{enumerate}

It may seem strange to put the chapter covering the planning process at the end
of the course, after students have learned the details of quantitative travel
modeling. The purpose for this is that I assign a term project where the
students build and calibrate a four-step model as they learn the techniques to do
so, and then complete an alternatives analysis using their models. To create
the time and space to do this project, we cover ``softer'' and conceptual topics
in the second half of the course.

The demonstration model the students calibrate and study is a model built in the
Cube travel modeling software for the Roanoke, Virginia, metropolitan region.
The model is a relatively advanced four-step, trip-based model with only 250
zones. The limited zone size means that the entire model system runs in
approximately 15 minutes on a laptop computer. I am grateful to Virginia DOT for
allowing my students the use of this model. Directions on how to use the Roanoke
model are given in the \protect\hyperlink{app-demomodel}{Appendices}.

A handful of assignments require the students to write numerical programs or
estimate statistical models. Some guidance on using R and RStudio to accomplish
these assignments is also given in the \protect\hyperlink{app-rstudio}{Appendices}.

\hypertarget{chap-blocks}{%
\chapter{Building Blocks}\label{chap-blocks}}

This chapter contains concepts, definitions, and mathematical techniques that will
be used throughout the semester.

\hypertarget{planning-for-human-systems}{%
\section{Planning for Human Systems}\label{planning-for-human-systems}}

If you look out on any sufficiently busy road, you will see a steady stream of
vehicles passing by. Each vehicle is largely indistinguishable from the others,
and it is easy as an engineer responsible for that road to see the cars driving
by as little more than an input to a problem. But the \emph{people} inside the cars
should not be indistinguishable from each other. Each person who is driving or
riding in each of those cars has their own reasons to be driving on that road.
One person might be driving to work; one person might be trying to get home to
his or her family. Another car might hold a family going on vacation, or a group
of friends heading to a movie.

If you don't recognize that each person who travels is
If the road is congested, then the
only apparent solution is to ``fix'' the road by expanding its capacity.

\hypertarget{travel-model-building-blocks}{%
\section{Travel Model Building Blocks}\label{travel-model-building-blocks}}

\hypertarget{travel-analysis-zones}{%
\subsection{Travel Analysis Zones}\label{travel-analysis-zones}}

SE Data

\hypertarget{highway-networks}{%
\subsection{Highway Networks}\label{highway-networks}}

Functional Type

Link capacity

Free-flow speed

Centroid connectors

\hypertarget{matrices}{%
\subsection{Matrices}\label{matrices}}

Skim matrices

OD matrices

\hypertarget{statistical-and-mathematical-techniques}{%
\section{Statistical and Mathematical Techniques}\label{statistical-and-mathematical-techniques}}

Many elements of travel modeling and forecasting require complex numerical and
quantitative techniques. In this section we will present some of these techniques.

\hypertarget{iterative-proportional-fitting}{%
\subsection{Iterative Proportional Fitting}\label{iterative-proportional-fitting}}

\hypertarget{regression-analysis}{%
\subsection{Regression Analysis}\label{regression-analysis}}

\hypertarget{numerical-optimization}{%
\subsection{Numerical Optimization}\label{numerical-optimization}}

\hypertarget{hw-blocks}{%
\section*{Homework}\label{hw-blocks}}
\addcontentsline{toc}{section}{Homework}

\begin{quote}
Some of these questions require a completed run of the demonstration model.
For instructions on accessing and running the model, see the \protect\hyperlink{app-demomodel}{Appendix}
\end{quote}

\begin{enumerate}
\def\labelenumi{\arabic{enumi}.}
\item
  With the TAZ layer and socioeconomic data in the demonstration model, make a
  set of choropleth maps showing: total households; household density; total jobs;
  job density; density of manufacturing vs office vs retail employment. Compare
  your maps with aerial imagery from Google Maps or OpenStreetMap. Describe the
  spatial patterns of the socioeconomic data in the model region. Identify which
  zones constitute the central business district, and identify any outlying
  employment centers.
\item
  With the highway network layer in the demonstration model, create maps
  showing: link functional type; link free flow speed; and link hourly capacity.
  Compare your maps with aerial imagery from Google Maps or OpenStreetMap.
  Identify the major freeways and principal arterials in the model region. \emph{Note}:
  you will need to run the demonstration model through the network setup step to
  calculate the capacities and append them to the link.
\item
  Find the shortest free-flow speed path along the network between two zones.
  Find the shortest distance path between the same two zones. Are the paths the
  same? Do the paths match what an online mapping service shows for a trip in the
  middle of the night?
\item
  Open the highway assignment report, which shows vehicle hours and miles
  traveled by facility type. What percent of the region's VMT occurs on freeways?
  What percent of the region's lane-miles are freeways?
\item
  Open the output highway network. Create a map of the
  highway links showing PM period level of service based on the volume to capacity
  ratios in the table below. How would you characterize traffic in Roanoke? Which
  is the worst-performing major facility?
\end{enumerate}

\hypertarget{chap-tripgen}{%
\chapter{Trip Generation}\label{chap-tripgen}}

\hypertarget{chap-distribution}{%
\chapter{Trip Distribution}\label{chap-distribution}}

\hypertarget{chap-modechoice}{%
\chapter{Mode and Destination Choice}\label{chap-modechoice}}

\hypertarget{chap-assignment}{%
\chapter{Network Assignment and Validation}\label{chap-assignment}}

\hypertarget{chap-process}{%
\chapter{The Planning Process}\label{chap-process}}

\hypertarget{appendix-appendix}{%
\appendix}


\hypertarget{app-demomodel}{%
\chapter{Demonstration Model}\label{app-demomodel}}

\hypertarget{running-the-model}{%
\section{Running the Model}\label{running-the-model}}

\hypertarget{files-and-reports}{%
\section{Files and Reports}\label{files-and-reports}}

\hypertarget{cube-tips-and-tricks}{%
\section{Cube Tips and Tricks}\label{cube-tips-and-tricks}}

\hypertarget{shortest-paths}{%
\subsection{Shortest Paths}\label{shortest-paths}}

\hypertarget{working-with-matrices}{%
\subsection{Working with Matrices}\label{working-with-matrices}}

\hypertarget{writing-custom-scripts}{%
\subsection{Writing Custom Scripts}\label{writing-custom-scripts}}

\hypertarget{app-rstudio}{%
\chapter{R and RStudio Help}\label{app-rstudio}}

R is a powerful, open-source statistical programming language used by both
professional and academic data scientists. It is among the computer languages
most suited to modern data science, and is growing rapidly in its user base and
available packages.

Some students may not feel comfortable working in a programming language like R or
a console-based application like RStudio, especially if they have used applications
primarily through a GUI.
This appendix provides a basic bootcamp for R and Rstudio, but cannot be a
comprehensive manual on RStudio, and it certainly cannot be one for R. Good
places to get more detailed help include:

\begin{itemize}
\tightlist
\item
  R help manuals
\item
  Stack Overflow
\end{itemize}

Some of the sections in this appendix are text-based, and some contain little
more than links to YouTube videos created by me or someone else.

\hypertarget{installing}{%
\section{Installing}\label{installing}}

There are two pieces of software you should install:

\begin{itemize}
\tightlist
\item
  \textbf{R} \url{https://cran.r-project.org/}: this contains
  the system libraries necessary to run R commands in a terminal on your computer,
  and contains a few additional helper applications. Install the most recent
  stable release for your operating system.
\item
  \textbf{RStudio} \url{https://rstudio.com/products/rstudio/download/} is an integrated application that makes using R considerably easier
  with text completion, file managment, and some GUI features.
\end{itemize}

Both software are available for Windows, MacOS, and Linux. The videos and screenshots
of the application I post will use MacOS; the R code for all systems is the same,
and the RStudio interface all systems is very similar with minor differences.

\hypertarget{rstudio-orientation}{%
\section{RStudio Orientation}\label{rstudio-orientation}}

\href{http://www.youtube.com/embed/c3xv8wOIj-g?rel=0}{\includegraphics{ceen565_files/figure-latex/rstudio-orientation-1.pdf}}

There is a very useful ``cheat sheet'' for working with RStudio on the \href{https://resources.rstudio.com/rstudio-cheatsheets/rstudio-ide-cheat-sheet}{RStudio}
website.

\hypertarget{r-packages}{%
\section{R Packages}\label{r-packages}}

One of the strengths of R is the ability for anyone to write packages. These
packages make it easier to read manipulate, and vizualize data; to estimate
statistical models; or to communicate results.

There are a number of ways to install additional packages. The most straightforward
is to use the \texttt{install.packages()} function in the console. The problems
in this book are solved with two additional packages\footnote{\texttt{tidyverse} is actually a collection
  of very useful packages, and many R users just load them all at once.}:

\begin{Shaded}
\begin{Highlighting}[]
\KeywordTok{install.packages}\NormalTok{(}\StringTok{"tidyverse"}\NormalTok{) }\CommentTok{# a suite of tools for data manipulation}
\KeywordTok{install.packages}\NormalTok{(}\StringTok{"mlogit"}\NormalTok{) }\CommentTok{# discrete choice modeling}
\end{Highlighting}
\end{Shaded}

RStudio also contains an interface to install and update packages.

Sometimes you want to use a package that has not yet been pushed to CRAN, the
international repository of ``approved'' R packages. This may be because the package
is in development, or for one reason or another does not meet CRAN's standards
for completeness, etc. Oftentimes, the package has been made available on GitHub.
You can install a package directly from GitHub with the \texttt{devtools} library. One
package you will want for the problems in the book is the \texttt{nhts2017} package
on the BYU Transportation GitHub account.

\begin{Shaded}
\begin{Highlighting}[]
\KeywordTok{install.packages}\NormalTok{(}\StringTok{"devtools"}\NormalTok{) }\CommentTok{# tools for installing development packages}
\NormalTok{devtools}\OperatorTok{::}\KeywordTok{install_github}\NormalTok{(}\StringTok{"gregmacfarlane/nhts2017"}\NormalTok{)}
\end{Highlighting}
\end{Shaded}

You only need to \textbf{install} a package once on your computer. But every time you
want to \textbf{use} a function in a package, you need to load the package with the
\texttt{library()} function. To load the \texttt{tidyverse} packages, for instance,

\begin{Shaded}
\begin{Highlighting}[]
\KeywordTok{library}\NormalTok{(tidyverse)}
\end{Highlighting}
\end{Shaded}

If you get errors when you run the command above, it means that for some reason
you did not install the package correctly. And if you ever get an error like

\begin{Shaded}
\begin{Highlighting}[]
\KeywordTok{kable}\NormalTok{(}\KeywordTok{tibble}\NormalTok{(}\DataTypeTok{x =} \DecValTok{1}\OperatorTok{:}\DecValTok{2}\NormalTok{, }\DataTypeTok{y =} \KeywordTok{c}\NormalTok{(}\StringTok{"blue"}\NormalTok{, }\StringTok{"red"}\NormalTok{)))}
\end{Highlighting}
\end{Shaded}

\begin{tabular}{r|l}
\hline
x & y\\
\hline
1 & blue\\
\hline
2 & red\\
\hline
\end{tabular}

It often means you didn't load the library. In this case, the \texttt{kable()} function
to make pretty tables is part of the \texttt{knitr} package.

\begin{Shaded}
\begin{Highlighting}[]
\KeywordTok{library}\NormalTok{(knitr)}
\KeywordTok{kable}\NormalTok{(}\KeywordTok{tibble}\NormalTok{(}\DataTypeTok{x =} \DecValTok{1}\OperatorTok{:}\DecValTok{2}\NormalTok{, }\DataTypeTok{y =} \KeywordTok{c}\NormalTok{(}\StringTok{"blue"}\NormalTok{, }\StringTok{"red"}\NormalTok{)))}
\end{Highlighting}
\end{Shaded}

\begin{tabular}{r|l}
\hline
x & y\\
\hline
1 & blue\\
\hline
2 & red\\
\hline
\end{tabular}

\hypertarget{working-with-tables}{%
\section{Working with Tables}\label{working-with-tables}}

Most data you will work with comes in a \emph{tabular} form, meaning that the data
is formatted in columns of variables and rows of observations.

\hypertarget{reading-data}{%
\subsection{Reading Data}\label{reading-data}}

Tabular data is often stored in a comma-separated values \texttt{.csv} file. To read a
data file like this in R, you can use the \texttt{read\_csv()} function included in
\texttt{tidyverse}.

\begin{Shaded}
\begin{Highlighting}[]
\NormalTok{trips <-}\StringTok{ }\KeywordTok{read_csv}\NormalTok{(}\StringTok{"data/demo_trips.csv"}\NormalTok{)}
\end{Highlighting}
\end{Shaded}

\begin{verbatim}
## Parsed with column specification:
## cols(
##   houseid = col_double(),
##   personid = col_character(),
##   trpmiles = col_double(),
##   trippurp = col_character()
## )
\end{verbatim}

\begin{Shaded}
\begin{Highlighting}[]
\KeywordTok{print}\NormalTok{(trips)}
\end{Highlighting}
\end{Shaded}

\begin{verbatim}
## # A tibble: 924 x 4
##     houseid personid trpmiles trippurp
##       <dbl> <chr>       <dbl> <chr>   
##  1 30182694 01         13.6   HBW     
##  2 40532989 02          9.38  HBSOCREC
##  3 40729475 01          5.06  HBSHOP  
##  4 40290784 01          0.509 HBSOCREC
##  5 30118876 02          0.599 NHB     
##  6 30352119 01         20.3   HBO     
##  7 30085077 01         22.5   NHB     
##  8 30180962 02          0.581 HBSOCREC
##  9 40155356 01          2.74  HBO     
## 10 30069734 01          4.65  NHB     
## # ... with 914 more rows
\end{verbatim}

This function will make a guess as to what the columns types should be. Often
we want to keep ID values as characters, even if they are numeric (this preserves
leading \texttt{0} values, etc.). We can tell \texttt{read\_csv()} what types we expect with
the \texttt{col\_types} argument.

\begin{Shaded}
\begin{Highlighting}[]
\NormalTok{trips <-}\StringTok{ }\KeywordTok{read_csv}\NormalTok{(}\StringTok{"data/demo_trips.csv"}\NormalTok{, }\DataTypeTok{col_types =} \KeywordTok{list}\NormalTok{(}\DataTypeTok{houseid =} \KeywordTok{col_character}\NormalTok{()))}
\KeywordTok{print}\NormalTok{(trips)}
\end{Highlighting}
\end{Shaded}

\begin{verbatim}
## # A tibble: 924 x 4
##    houseid  personid trpmiles trippurp
##    <chr>    <chr>       <dbl> <chr>   
##  1 30182694 01         13.6   HBW     
##  2 40532989 02          9.38  HBSOCREC
##  3 40729475 01          5.06  HBSHOP  
##  4 40290784 01          0.509 HBSOCREC
##  5 30118876 02          0.599 NHB     
##  6 30352119 01         20.3   HBO     
##  7 30085077 01         22.5   NHB     
##  8 30180962 02          0.581 HBSOCREC
##  9 40155356 01          2.74  HBO     
## 10 30069734 01          4.65  NHB     
## # ... with 914 more rows
\end{verbatim}

You can also write tables back to \texttt{.csv} with the \texttt{write\_csv()} command.

\hypertarget{modifying-and-summarizing-tables}{%
\subsection{Modifying and Summarizing Tables}\label{modifying-and-summarizing-tables}}

In much of this section, we will work with the \texttt{nhts\_trips} dataset of trips
from the 2017 National Household Travel Survey in the \texttt{nhts2017} package you
installed from GitHub above.

\begin{Shaded}
\begin{Highlighting}[]
\KeywordTok{library}\NormalTok{(nhts2017)}
\NormalTok{trips <-}\StringTok{ }\NormalTok{nhts_trips}
\NormalTok{trips}
\end{Highlighting}
\end{Shaded}

\begin{verbatim}
## # A tibble: 923,572 x 62
##    houseid personid tdtrpnum strttime            endtime            
##    <chr>   <chr>       <dbl> <dttm>              <dttm>             
##  1 300000~ 01              1 2017-10-10 10:00:00 2017-10-10 10:15:00
##  2 300000~ 01              2 2017-10-10 15:10:00 2017-10-10 15:30:00
##  3 300000~ 02              1 2017-10-10 07:00:00 2017-10-10 09:00:00
##  4 300000~ 02              2 2017-10-10 18:00:00 2017-10-10 20:30:00
##  5 300000~ 03              1 2017-10-10 08:45:00 2017-10-10 09:00:00
##  6 300000~ 03              2 2017-10-10 14:30:00 2017-10-10 14:45:00
##  7 300000~ 01              1 2017-10-10 11:15:00 2017-10-10 11:30:00
##  8 300000~ 01              2 2017-10-10 23:30:00 2017-10-10 23:40:00
##  9 300000~ 01              1 2017-10-10 05:50:00 2017-10-10 06:05:00
## 10 300000~ 01              2 2017-10-10 07:00:00 2017-10-10 07:15:00
## # ... with 923,562 more rows, and 57 more variables: trvlcmin <dbl+lbl>,
## #   trpmiles <dbl+lbl>, trptrans <chr+lbl>, trpaccmp <dbl+lbl>,
## #   trphhacc <dbl+lbl>, vehid <chr+lbl>, trwaittm <dbl+lbl>,
## #   numtrans <dbl+lbl>, tracctm <dbl+lbl>, drop_prk <chr+lbl>,
## #   tregrtm <dbl+lbl>, whodrove <chr+lbl>, whyfrom <chr+lbl>,
## #   loop_trip <chr+lbl>, trphhveh <chr+lbl>, hhmemdrv <chr+lbl>,
## #   hh_ontd <dbl+lbl>, nonhhcnt <dbl+lbl>, numontrp <dbl+lbl>,
## #   psgr_flg <chr+lbl>, pubtrans <chr+lbl>, trippurp <chr+lbl>,
## #   dweltime <dbl+lbl>, tdwknd <chr+lbl>, vmt_mile <dbl+lbl>,
## #   drvr_flg <chr+lbl>, whytrp1s <chr+lbl>, ontd_p1 <chr+lbl>,
## #   ontd_p2 <chr+lbl>, ontd_p3 <chr+lbl>, ontd_p4 <chr+lbl>,
## #   ontd_p5 <chr+lbl>, ontd_p6 <chr+lbl>, ontd_p7 <chr+lbl>,
## #   ontd_p8 <chr+lbl>, ontd_p9 <chr+lbl>, ontd_p10 <chr+lbl>,
## #   ontd_p11 <chr+lbl>, ontd_p12 <chr+lbl>, ontd_p13 <chr+lbl>,
## #   tdcaseid <chr>, tracc_wlk <chr+lbl>, tracc_pov <chr+lbl>,
## #   tracc_bus <chr+lbl>, tracc_crl <chr+lbl>, tracc_sub <chr+lbl>,
## #   tracc_oth <chr+lbl>, tregr_wlk <chr+lbl>, tregr_pov <chr+lbl>,
## #   tregr_bus <chr+lbl>, tregr_crl <chr+lbl>, tregr_sub <chr+lbl>,
## #   tregr_oth <chr+lbl>, whyto <chr+lbl>, gasprice <chr>, wttrdfin <dbl>,
## #   whytrp90 <chr+lbl>
\end{verbatim}

\hypertarget{select-filter-and-chains}{%
\subsubsection{Select, Filter, and Chains}\label{select-filter-and-chains}}

This table is pretty overwhelming. But there are two functions that can help
us pare it down:

\begin{itemize}
\tightlist
\item
  \texttt{select()} lets you select columns in a table using the names of the columns.
\item
  \texttt{filter()} lets you select rows in a table that meet a certain condition.
\end{itemize}

Let's practice this by selecting our \texttt{trips} dataset to only include the id
columns, the trip length, and the trip purpose.

\begin{Shaded}
\begin{Highlighting}[]
\KeywordTok{select}\NormalTok{(trips, houseid, personid, trpmiles, trippurp)}
\end{Highlighting}
\end{Shaded}

\begin{verbatim}
## # A tibble: 923,572 x 4
##    houseid  personid  trpmiles trippurp                                    
##    <chr>    <chr>    <dbl+lbl> <chr+lbl>                                   
##  1 30000007 01            5.24 HBO [Home-based trip (other)]               
##  2 30000007 01            5.15 HBO [Home-based trip (other)]               
##  3 30000007 02           84.0  HBW [Home-based trip (work)]                
##  4 30000007 02           81.6  HBW [Home-based trip (work)]                
##  5 30000007 03            2.25 HBO [Home-based trip (other)]               
##  6 30000007 03            2.24 HBO [Home-based trip (other)]               
##  7 30000008 01            8.02 HBW [Home-based trip (work)]                
##  8 30000008 01            8.02 HBW [Home-based trip (work)]                
##  9 30000012 01            3.40 HBSOCREC [Home-based trip (social/recreatio~
## 10 30000012 01            3.40 HBSOCREC [Home-based trip (social/recreatio~
## # ... with 923,562 more rows
\end{verbatim}

Let's also practice filtering the \texttt{trips} dataset to only include trips
of the purpose ``HBO'' (home-based other). Notice how the number of rows
in the table trips is much smaller.

\begin{Shaded}
\begin{Highlighting}[]
\KeywordTok{filter}\NormalTok{(trips, trippurp }\OperatorTok{==}\StringTok{ "HBW"}\NormalTok{) }\CommentTok{# use double equals as comparison}
\end{Highlighting}
\end{Shaded}

\begin{verbatim}
## # A tibble: 117,368 x 62
##    houseid personid tdtrpnum strttime            endtime            
##    <chr>   <chr>       <dbl> <dttm>              <dttm>             
##  1 300000~ 02              1 2017-10-10 07:00:00 2017-10-10 09:00:00
##  2 300000~ 02              2 2017-10-10 18:00:00 2017-10-10 20:30:00
##  3 300000~ 01              1 2017-10-10 11:15:00 2017-10-10 11:30:00
##  4 300000~ 01              2 2017-10-10 23:30:00 2017-10-10 23:40:00
##  5 300000~ 01              5 2017-10-10 09:00:00 2017-10-10 09:20:00
##  6 300000~ 01              7 2017-10-10 15:30:00 2017-10-10 16:05:00
##  7 300000~ 01              1 2017-10-10 08:00:00 2017-10-10 08:20:00
##  8 300000~ 01              2 2017-10-10 18:00:00 2017-10-10 20:00:00
##  9 300000~ 02              3 2017-10-10 09:00:00 2017-10-10 11:00:00
## 10 300000~ 02              4 2017-10-10 18:30:00 2017-10-10 20:30:00
## # ... with 117,358 more rows, and 57 more variables: trvlcmin <dbl+lbl>,
## #   trpmiles <dbl+lbl>, trptrans <chr+lbl>, trpaccmp <dbl+lbl>,
## #   trphhacc <dbl+lbl>, vehid <chr+lbl>, trwaittm <dbl+lbl>,
## #   numtrans <dbl+lbl>, tracctm <dbl+lbl>, drop_prk <chr+lbl>,
## #   tregrtm <dbl+lbl>, whodrove <chr+lbl>, whyfrom <chr+lbl>,
## #   loop_trip <chr+lbl>, trphhveh <chr+lbl>, hhmemdrv <chr+lbl>,
## #   hh_ontd <dbl+lbl>, nonhhcnt <dbl+lbl>, numontrp <dbl+lbl>,
## #   psgr_flg <chr+lbl>, pubtrans <chr+lbl>, trippurp <chr+lbl>,
## #   dweltime <dbl+lbl>, tdwknd <chr+lbl>, vmt_mile <dbl+lbl>,
## #   drvr_flg <chr+lbl>, whytrp1s <chr+lbl>, ontd_p1 <chr+lbl>,
## #   ontd_p2 <chr+lbl>, ontd_p3 <chr+lbl>, ontd_p4 <chr+lbl>,
## #   ontd_p5 <chr+lbl>, ontd_p6 <chr+lbl>, ontd_p7 <chr+lbl>,
## #   ontd_p8 <chr+lbl>, ontd_p9 <chr+lbl>, ontd_p10 <chr+lbl>,
## #   ontd_p11 <chr+lbl>, ontd_p12 <chr+lbl>, ontd_p13 <chr+lbl>,
## #   tdcaseid <chr>, tracc_wlk <chr+lbl>, tracc_pov <chr+lbl>,
## #   tracc_bus <chr+lbl>, tracc_crl <chr+lbl>, tracc_sub <chr+lbl>,
## #   tracc_oth <chr+lbl>, tregr_wlk <chr+lbl>, tregr_pov <chr+lbl>,
## #   tregr_bus <chr+lbl>, tregr_crl <chr+lbl>, tregr_sub <chr+lbl>,
## #   tregr_oth <chr+lbl>, whyto <chr+lbl>, gasprice <chr>, wttrdfin <dbl>,
## #   whytrp90 <chr+lbl>
\end{verbatim}

One \textbf{\emph{extremely}} useful feature of the \texttt{tidyverse} functions is the chain
operator, \texttt{\%\textgreater{}\%}. This operator basically does the opposite of the assigment
operator \texttt{\textless{}-}. While assignment says ``take the thing on the right and put it in
the thing on the left,'' chain says ``take the thing on the left and pass it as
the first argument of the function on the right.'' What this means in practice is
we can chain R commands together. So we can do the \texttt{select} \emph{and} the \texttt{filter}
statements in sequence,

\begin{Shaded}
\begin{Highlighting}[]
\NormalTok{trips }\OperatorTok{%>%}
\StringTok{  }\KeywordTok{select}\NormalTok{(houseid, personid, trpmiles, trippurp) }\OperatorTok{%>%}
\StringTok{  }\KeywordTok{filter}\NormalTok{(trippurp }\OperatorTok{==}\StringTok{ "HBW"}\NormalTok{)}
\end{Highlighting}
\end{Shaded}

\begin{verbatim}
## # A tibble: 117,368 x 4
##    houseid  personid  trpmiles trippurp                    
##    <chr>    <chr>    <dbl+lbl> <chr+lbl>                   
##  1 30000007 02           84.0  HBW [Home-based trip (work)]
##  2 30000007 02           81.6  HBW [Home-based trip (work)]
##  3 30000008 01            8.02 HBW [Home-based trip (work)]
##  4 30000008 01            8.02 HBW [Home-based trip (work)]
##  5 30000012 01            4.29 HBW [Home-based trip (work)]
##  6 30000012 01            6.82 HBW [Home-based trip (work)]
##  7 30000039 01           11.5  HBW [Home-based trip (work)]
##  8 30000041 01           73.7  HBW [Home-based trip (work)]
##  9 30000041 02           77.9  HBW [Home-based trip (work)]
## 10 30000041 02           77.8  HBW [Home-based trip (work)]
## # ... with 117,358 more rows
\end{verbatim}

Notice that we didn't have to tell the \texttt{select} and \texttt{filter} functions the
name of the table we were selecting or filtering. The \texttt{\%\textgreater{}\%} chain operator did
that for us.

Once we have the table we want, we can assign it to a new object called \texttt{mytrips}
In this case, let's get \texttt{HBO} and \texttt{HBW} trips.

\begin{Shaded}
\begin{Highlighting}[]
\NormalTok{mytrips <-}\StringTok{ }\NormalTok{trips }\OperatorTok{%>%}
\StringTok{  }\KeywordTok{select}\NormalTok{(houseid, personid, trpmiles, trippurp) }\OperatorTok{%>%}
\StringTok{  }\KeywordTok{filter}\NormalTok{(trippurp }\OperatorTok{%in%}\StringTok{ }\KeywordTok{c}\NormalTok{(}\StringTok{"HBW"}\NormalTok{, }\StringTok{"HBO"}\NormalTok{)) }\CommentTok{# use %in% for multiple comparisons.}
\end{Highlighting}
\end{Shaded}

\hypertarget{mutate-summarize-and-group}{%
\subsection{Mutate, Summarize, and Group}\label{mutate-summarize-and-group}}

Sometimes we want to calculate a new column in a table, or recompute an
existing column. We can do that with the \texttt{mutate} function, and we can
put more than one calculation in a single \texttt{mutate} statement.

\begin{Shaded}
\begin{Highlighting}[]
\NormalTok{mytrips }\OperatorTok{%>%}
\StringTok{  }\KeywordTok{mutate}\NormalTok{(}
    \DataTypeTok{tripkm =}\NormalTok{ trpmiles }\OperatorTok{*}\StringTok{ }\FloatTok{1.60934}\NormalTok{, }\CommentTok{# convert miles to km.}
    \DataTypeTok{longtrip =} \KeywordTok{ifelse}\NormalTok{(tripkm }\OperatorTok{>}\StringTok{ }\DecValTok{50}\NormalTok{, }\OtherTok{TRUE}\NormalTok{, }\OtherTok{FALSE}\NormalTok{) }\CommentTok{# is trip longer than 50 km?}
\NormalTok{  )}
\end{Highlighting}
\end{Shaded}

\begin{verbatim}
## # A tibble: 307,390 x 6
##    houseid  personid  trpmiles trippurp                     tripkm longtrip
##    <chr>    <chr>    <dbl+lbl> <chr+lbl>                  <dbl+lb> <lgl>   
##  1 30000007 01            5.24 HBO [Home-based trip (oth~     8.44 FALSE   
##  2 30000007 01            5.15 HBO [Home-based trip (oth~     8.29 FALSE   
##  3 30000007 02           84.0  HBW [Home-based trip (wor~   135.   TRUE    
##  4 30000007 02           81.6  HBW [Home-based trip (wor~   131.   TRUE    
##  5 30000007 03            2.25 HBO [Home-based trip (oth~     3.62 FALSE   
##  6 30000007 03            2.24 HBO [Home-based trip (oth~     3.61 FALSE   
##  7 30000008 01            8.02 HBW [Home-based trip (wor~    12.9  FALSE   
##  8 30000008 01            8.02 HBW [Home-based trip (wor~    12.9  FALSE   
##  9 30000012 01            4.29 HBW [Home-based trip (wor~     6.91 FALSE   
## 10 30000012 01            6.82 HBW [Home-based trip (wor~    11.0  FALSE   
## # ... with 307,380 more rows
\end{verbatim}

Other times we want to calculate summary statistics like means.
For this we can use the \texttt{summarize()} function.

\begin{Shaded}
\begin{Highlighting}[]
\NormalTok{mytrips }\OperatorTok{%>%}
\StringTok{  }\KeywordTok{summarize}\NormalTok{(}
    \DataTypeTok{mean_trip =} \KeywordTok{mean}\NormalTok{(trpmiles),}
    \DataTypeTok{sd_trip =} \KeywordTok{sd}\NormalTok{(trpmiles),}
    \DataTypeTok{max_trip =} \KeywordTok{max}\NormalTok{(trpmiles),}
    \DataTypeTok{min_trip =} \KeywordTok{min}\NormalTok{(trpmiles)}
\NormalTok{  )}
\end{Highlighting}
\end{Shaded}

\begin{verbatim}
## # A tibble: 1 x 4
##   mean_trip sd_trip max_trip min_trip
##       <dbl>   <dbl>    <dbl>    <dbl>
## 1      9.81    32.1    5699.       -9
\end{verbatim}

Finally, we sometimes want to calculate summary statistics for different groups.
We can tell \texttt{tidyverse} to group our tables with the \texttt{group\_by()} function.

\begin{Shaded}
\begin{Highlighting}[]
\NormalTok{mytrips }\OperatorTok{%>%}
\StringTok{  }\KeywordTok{group_by}\NormalTok{(trippurp) }\OperatorTok{%>%}
\StringTok{  }\KeywordTok{summarize}\NormalTok{(}
    \DataTypeTok{mean_trip =} \KeywordTok{mean}\NormalTok{(trpmiles),}
    \DataTypeTok{sd_trip =} \KeywordTok{sd}\NormalTok{(trpmiles),}
    \DataTypeTok{max_trip =} \KeywordTok{max}\NormalTok{(trpmiles),}
    \DataTypeTok{min_trip =} \KeywordTok{min}\NormalTok{(trpmiles)}
\NormalTok{  )}
\end{Highlighting}
\end{Shaded}

\begin{verbatim}
## # A tibble: 2 x 5
##   trippurp                      mean_trip sd_trip max_trip min_trip
## * <chr+lbl>                         <dbl>   <dbl>    <dbl>    <dbl>
## 1 HBO [Home-based trip (other)]      7.73    32.0    5699.       -9
## 2 HBW [Home-based trip (work)]      13.2     31.9    2927.       -9
\end{verbatim}

\begin{quote}
As you might expect, work trips are on average longer than other kinds of trips.
But some people report very long trips! You might want to filter your data more
carefully for real analyses.
\end{quote}

\hypertarget{graphics-with-ggplot}{%
\section{\texorpdfstring{Graphics with \texttt{ggplot}}{Graphics with ggplot}}\label{graphics-with-ggplot}}

\bibliography{book.bib}


\end{document}
